We provide a proof of the claim \ref{clm:sc-deg-src}.
\begin{proof}
    Given $\alpha\textit{S-SourceOrExcess}(\beta,1)$, in order to find
    the minimum number of solutions, we can assume that we there are no other sources other than the first $\alpha$ nodes.
    We argue that the minimum number of solutions is $\mu = \lceil \alpha/\beta \rceil$ since we can create 
    $\mu$ sinks that connect to at most $\beta$ sources each. We will refer to these
    graphs as sink graphs.

    Furthermore, we can demonstrate that any graph can be reduced to a sink graph.
    Our first observation is to convert every excess node to a sink node. To do that we can apply two simplification 
    phases:
    \begin{enumerate}
        \item Removal phase: Remove all paths between two excess nodes or between excess and sink nodes.
        We can observe that after this steps our graph only contains sources and sinks.
        \item Merging phase: If there are sinks with degrees $(a,0)$ and $(b,0)$ such that
        $a + b \leq \beta$ then we can merge the two to remove a solution. Keep repeating this process until
        no two sinks can be merged
    \end{enumerate}
    We can observe that any graph will simplify to a sink graph, and therefore we will always have at least $\mu$ solutions,
    as each step will never increase the solution set.
\end{proof}